\chapter{Le théorème de Pythagore}
\label{corriges:13}
\section*{Exercice 13.1 -- Une diagonale}
La diagonale forme l'hypoténuse d'un triangle rectangle de côtés $6$ et $11$. Elle vaut $\sqrt{6^2+11^2}=\sqrt{157}\unite{cm}\approx12{,}52996\unite{cm}$. Au millimètre près : $12{,}5\unite{cm}$, soit $125\unite{mm}$.

\section*{Exercice 13.2 -- Une échelle}
Le triangle est rectangle. À $5\unite m$, $h^2=13^2-5^2=144$, donc $h=12\unite m$. À $8\unite m$, $h=\sqrt{169-64}=\sqrt{105}\unite m\approx10{,}25\unite m$. Éloigner le pied diminue ici la hauteur atteinte ; le contrôle du sens confirme le calcul.

\section*{Exercice 13.3 -- Rectangle ou non ?}
Pour $(6,8,10)$, $10^2=100=6^2+8^2$ : par la réciproque de Pythagore, le triangle est rectangle en face du côté $10$. Pour $(6,8,11)$, $11^2=121\ne100$ : par contraposée du théorème direct, il n'est pas rectangle. Les trois longueurs de chaque liste respectent l'inégalité triangulaire.

\section*{Exercice 13.4 -- Distance dans un repère}
Les écarts horizontal et vertical valent $6$ et $8$. Pythagore donne $AB=\sqrt{6^2+8^2}=10$ unités. Le repère doit avoir des axes perpendiculaires et une même unité sur chacun, sans quoi les écarts de coordonnées ne représentent pas directement les deux longueurs du triangle.
